\documentclass[a4paper,12pt]{article}
\usepackage[utf8]{inputenc}
\usepackage[T1]{fontenc}
\usepackage{amsmath,amssymb}
\usepackage{xcolor}
\usepackage{geometry}
\geometry{margin=2.5cm}

\title{Resumo de Equações Diferenciais}
\author{}
\date{}

\begin{document}
\maketitle

\section*{Classificação das Equações Diferenciais}
\begin{itemize}
    \item \textbf{Ordem} $\to$ Maior derivada presente
    \item \textbf{Grau} $\to$ expoente da maior derivada na forma normal
    \item \textbf{Linearidade} $\to$ Uma ED é \textbf{linear} se $y$ e derivadas aparecem em potência 1, sem produtos entre si.
\end{itemize}

Forma geral linear de ordem $n$:
\[
a_n(x)y^{(n)}+a_{n-1}(x)y^{(n-1)}+\dots+a_1(x)y'+a_0(x)y=g(x)
\]

\section*{Equações de Primeira Ordem}

\subsection*{Separáveis (não lineares)}
Forma:
\[
\frac{dy}{dt}=g(t)h(y) \quad \to \quad \frac{dy}{h(y)}=g(t)\,dt
\]

Solução:
\[
\ln|y|= \int g(t)\,dt + c
\]
\[
\textcolor{purple}{y(t)=c\cdot e^{\int g(t)\,dt}}
\]

\subsection*{Lineares}
Forma:
\[
y'+a(t)y=b(t)
\]

\textbf{Fator integrante:}
\[
\textcolor{blue}{\mu(t)=e^{\int a(t)\,dt}}
\]

Solução:
\[
\textcolor{purple}{y\mu=\int \mu\, b(t)\,dt + C}
\]

\subsection*{Exatas}
Forma:
\[
M(x,y)+N(x,y)\,y'=0
\]

Verificar:
\[
\textcolor{blue}{\frac{\partial M}{\partial y}=\frac{\partial N}{\partial x}}
\]

Se sim, é exata. Existe função $\psi(x,y)$ tal que:
\[
d\psi = M\,dx+N\,dy
\]

Solução:
\[
\textcolor{purple}{\psi(x,y)=C}
\]

Resolução:
\[
N(x,y) = \frac{\partial}{\partial y} \int M(x,y)\,dx + \frac{\partial}{\partial y} C(y)
\]

\section*{Equações de Segunda Ordem}
Forma geral:
\[
y''+a\,y'+b\,y=f(x)
\]

\subsection*{Solução Homogênea}
\[
y''+a\,y'+b\,y=0
\]

Equação característica:
\[
\textcolor{blue}{r^2+a\,r+b=0}
\]

\begin{itemize}
    \item \textbf{Raízes reais diferentes:}
    \[
    \textcolor{purple}{y(t)=c_1e^{r_1t}+c_2e^{r_2t}}
    \]
    \item \textbf{Raízes reais iguais:}
    \[
    \textcolor{purple}{y(t)=(c_1+c_2t)e^{rt}}
    \]
    \item \textbf{Raízes Complexas:} $r_{1,2}=\lambda \pm \mu i$
    \[
    \textcolor{purple}{y(t)=e^{\lambda t}(c_1\cos \mu t + c_2\sin \mu t)}
    \]
\end{itemize}

\subsection*{Solução Não-homogênea}
Forma:
\[
y(t) = Y_h(t) + Y_P(t)
\]

\subsubsection*{Método dos Coeficientes Indeterminados}
\[
a\,y''+b\,y'+c\,y=g(t)
\]

\paragraph{Tabela de Hipóteses}
\begin{center}
\begin{tabular}{c|c}
$g(t)$ & $Y(t)$ \\ \hline
$k$ & $A$ \\
$e^{at}$ & $Ae^{at}$ \\
$t e^{at}$ & $(At+B)e^{at}$ \\
$\sin \beta t$ ou $\cos \beta t$ & $A\cos \beta t + B\sin \beta t$ \\
$t^n$ (polinômio) & polinômio mesmo grau $n$
\end{tabular}
\end{center}

Se repetir termo da homogênea: multiplicar por $t$.

\subsubsection*{Método da Variação de Parâmetros}
Para:
\[
y''+p(t)y'+q(t)y=g(t)
\]

Se $y_1,y_2$ são soluções homogêneas:
\[
\textcolor{blue}{Y_p=u_1y_1+u_2y_2}
\]
onde
\[
\textcolor{blue}{u_1'=\frac{-g(t) \cdot y_2}{W}}, \qquad 
\textcolor{blue}{u_2'=\frac{g(t) \cdot y_1}{W}}
\]
Wronskiano:
\[
W=y_1y_2'-y_2y_1'
\]

Solução:
\[
\textcolor{purple}{Y_p =\int\frac{-g(t) \cdot y_2}{W}\,dt \cdot y_1 +\int\frac{g(t) \cdot y_1}{W}\,dt \cdot y_2}
\]

\section*{Oscilador Harmônico}
Modelo:
\[
m x'' + c x' + k x = F(t)
\]
onde:
\begin{itemize}
    \item $m$ $\to$ massa
    \item $c$ $\to$ amortecimento
    \item $k$ $\to$ mola
    \item $F(t)$ $\to$ força externa
\end{itemize}

\subsection*{Caso Livre sem amortecimento}
\[
m x''+k x=0 \quad\Rightarrow\quad x''+\omega^2 x=0,\qquad \omega=\sqrt{\frac{k}{m}}
\]
Solução:
\[
x=C_1\cos(\omega t)+C_2\sin(\omega t)
\]

\subsection*{Caso Amortecido}
\[
m x''+c x'+k x=0
\]
Equação característica:
\[
m r^2+c r+k=0 \quad\Rightarrow\quad r_{1,2}=\frac{-c \pm \sqrt{c^2-4km}}{2m}
\]
\begin{center}
\begin{tabular}{|c|c|c|c|}
\hline
\textbf{Caso} & $\Delta$ & $x(t)$ & \textbf{Comportamento} \\
\hline
Superamortecido & $c^2>4mk$ & $x(t)=Ae^{r_1t}+Be^{r_2t}$ & Queda exponencial sem oscilar \\
\hline
Criticamente amortecido & $c^2=4mk$ & $x(t)=Ae^{\frac{-c}{2m}t}+t B e^{\frac{-c}{2m}t}$ & Queda após pico sem oscilar \\
\hline
Subamortecido & $c^2<4mk$ & $\begin{array}{l}
x(t)=e^{\frac{-c}{2m}t}\Bigl(A\cos\bigl(\frac{\sqrt{4km-c^2}}{2m}t\bigr)\\[2mm]
\quad + B\sin\bigl(\frac{\sqrt{4km-c^2}}{2m}t\bigr)\Bigr)
\end{array}$ & Oscila com amplitude decrescente. \\
\hline
\end{tabular}
\end{center}

\subsection*{Ressonância}
\[
m x''+c x'+k x=F_0\cos(\omega t)
\]
Ocorre quando frequência externa = natural.

\end{document}